\section{Описание датасетов}\label{app:A}

\subsection{DEAM (Database for Emotional Analysis of Music)}

\begin{table}[H]
\centering
\caption{Характеристики DEAM}
\begin{tabular}{L{6cm}L{8cm}}
\toprule
Характеристика & Значение \\
\midrule
Источник & MediaEval 2013–2018 challenges, \url{http://cvml.unige.ch/databases/DEAM} \\
Лицензия & Creative Commons (треки из Free Music Archive) \\
Размер & 1802 фрагмента $\times$ 45 секунд каждый \\
Sample rate & 44100\,Гц (мы передискретизируем до 22050\,Гц) \\
Аннотации & Покадровые непрерывные arousal/valence ($\sim$ 60 фреймов на трек, шаг 500\,мс) \\
Inter-annotator agreement (Pearson $\rho$) & arousal: 0.59, valence: 0.55 \\
Распределение жанров (примерно) & Pop $\sim$25\%, Rock $\sim$20\%, Electronic $\sim$15\%, Classical $\sim$10\%, Hip-hop $\sim$8\%, Other $\sim$22\% \\
\bottomrule
\end{tabular}
\end{table}

\textbf{Применение в работе.} Используется для обучения голов arousal\_cls (3 класса), valence\_cls (3 класса), arousal\_reg (regression), valence\_reg (regression). Дискретизация порогами $\pm 0.2$.

\subsection{Harmonix Set}

\begin{table}[H]
\centering
\caption{Характеристики Harmonix Set}
\begin{tabular}{L{6cm}L{8cm}}
\toprule
Характеристика & Значение \\
\midrule
Источник & Nieto et al., 2019; \url{https://github.com/urinieto/harmonixset} \\
Лицензия & Аннотации — MIT, аудио — отсутствует (требуется загрузка через YouTube) \\
Размер исходный & 912 треков с структурной разметкой \\
Размер использован & 743 трека (после загрузки через yt-dlp; 169 удалены с YouTube) \\
Длительность среднего трека & 3:43 \\
Аннотации & Beat positions, downbeats, section labels (17 категорий → сжаты до 6) \\
Inter-annotator agreement & section labels: $\sim$75\% \\
\bottomrule
\end{tabular}
\end{table}

\textbf{Маппинг 17 → 6 классов:}

\begin{lstlisting}
intro, intro_a, intro_b               -> Intro
verse, verse_a, verse_b, theme        -> Verse
bridge, transition, transition_*      -> Bridge
chorus, chorus_a, chorus_b, hook      -> Chorus
instrumental, solo, breakdown         -> Instrumental
outro, fadeout                        -> Outro
\end{lstlisting}

\subsection{GTZAN}

\begin{table}[H]
\centering
\caption{Характеристики GTZAN}
\begin{tabular}{L{6cm}L{8cm}}
\toprule
Характеристика & Значение \\
\midrule
Источник & Tzanetakis \& Cook, 2002; \url{http://marsyas.info/downloads/datasets.html} \\
Лицензия & Public domain (academic use) \\
Размер & 1000 файлов wav, 30 секунд каждый \\
Sample rate & 22050\,Гц (mono) \\
Жанры & blues, classical, country, disco, hiphop, jazz, metal, pop, reggae, rock (по 100 треков) \\
Использовано & 999 (один файл повреждён) \\
Известные дубликаты & $\sim$5\% \\
\bottomrule
\end{tabular}
\end{table}

\section{Гиперпараметры экспериментов}\label{app:B}

\subsection{Финальная конфигурация (default.yaml для CNN/PANNs/BiLSTM)}

\begin{lstlisting}
audio:
  sr: 22050
  mono: true

features:
  mode: "log_mel"
  n_fft: 2048
  hop_length: 512
  win_length: 2048
  n_mels: 128
  fmin: 30.0
  fmax: null

windowing:
  window_seconds: 10.0
  hop_seconds: 1.0

model:
  in_channels: 1
  n_segment_classes: 6
  n_arousal_classes: 3
  n_valence_classes: 3
  n_genre_classes: 10
  dropout: 0.4
  backbone_dropout: 0.25
  enable_regression: true

training:
  seed: 42
  lr: 1.0e-3
  weight_decay: 5.0e-4
  batch_size: 64
  epochs: 50
  early_stopping_patience: 15
  scheduler: cosine_warm_restarts
  loss_type: focal
  focal_gamma: 1.5
  mixup_alpha: 0.2
  max_class_weight: 2.5

spec_augment:
  freq_mask: 27
  time_mask: 60
  n_masks: 2

loss_weights:
  segment: 1.5
  arousal_cls: 1.0
  valence_cls: 1.0
  arousal_reg: 0.5
  valence_reg: 0.5
  genre: 1.0

postprocess:
  smooth_kernel: 5
  min_segment_duration: 4.0
  use_viterbi: true
  transition_penalty: 1.0
  use_position_priors: true
  smooth_kernel_segment: 5
  min_segment_duration_segment: 6.0
\end{lstlisting}

\subsection{AST v2 (ast.yaml)}

Отличается от default следующим:

\begin{lstlisting}
architecture: "ast"

model:
  pretrained_model: "MIT/ast-finetuned-audioset-10-10-0.4593"
  freeze_layers: 8
  hidden_dim: 768
  shared_fc_dim: 256
  dropout: 0.3

training:
  lr: 5.0e-5         # backbone
  lr_head: 1.0e-3    # heads
  batch_size: 16
  grad_accumulation_steps: 4
  epochs: 30
  early_stopping_patience: 10
  weight_decay: 0.01
\end{lstlisting}

\subsection{CNN v3 (failed experiment, cnn\_v3.yaml)}

Отличается от default только следующим:

\begin{lstlisting}
training:
  focal_gamma: 2.5            # 1.5 in default
  max_class_weight: 3.5       # 2.5 in default

loss_weights:
  segment: 3.5                # 1.5 in default
\end{lstlisting}

\section{Confusion matrices}\label{app:C}

См. графические результаты в \code{results/plots/}:

\begin{itemize}
  \item \code{confusion\_segment\_cnn\_v2.png} — CNN v2, сегментная задача
  \item \code{confusion\_segment\_ast\_v2.png} — AST v2, сегментная задача (для сравнения)
  \item \code{confusion\_arousal\_v2.png}, \code{confusion\_valence\_v2.png}, \code{confusion\_genre\_v2.png} — для CNN v2
  \item \code{comparison\_4tasks.png} — сравнение 6 моделей по 4 задачам в виде гистограммы
  \item \code{delta\_v1\_v2.png} — изменения метрик v1 → v2
\end{itemize}

\section{Скриншоты прототипа}\label{app:D}

Снимки экрана веб-прототипа (можно делать через preview\_screenshot или браузерную консоль):

\begin{enumerate}
  \item \textbf{Главная страница} — список треков с поиском, кнопкой загрузки.
  \item \textbf{Dashboard трека} (Dynamite) — все 6 панелей: структура, эмоции, AV-trajectory, жанр, спектрограмма, сводка.
  \item \textbf{Top-3 жанров с warning} — пример low-confidence классификации с жёлтой плашкой.
  \item \textbf{AV-trajectory с кометой} — Russell circumplex с движущейся точкой и шлейфом за последние 20 секунд.
  \item \textbf{Mel-spectrogram} — viridis-окрашенная спектрограмма на 3-минутном треке.
  \item \textbf{Loading skeletons} — пример отображения skeleton-screens во время загрузки.
\end{enumerate}

\section{JSON-схема timeline}\label{app:E}

Структура ответа \code{/api/tracks/:id}:

\begin{lstlisting}
{
  "id": "uuid-string",
  "filename": "uuid-string.mp3",
  "originalName": "0375_dynamite.mp3",
  "status": "ready",
  "createdAt": "2026-04-15T...",
  "timeline": {
    "metadata": {
      "duration_sec": 204.1,
      "window_seconds": 10.0,
      "hop_seconds": 1.0
    },
    "segment": [
      { "start": 0.0, "end": 8.0, "label": "Chorus", "confidence": 0.63 },
      { "start": 8.0, "end": 24.0, "label": "Verse", "confidence": 0.54 },
      ...
    ],
    "arousal": [ /* same shape */ ],
    "valence": [ /* same shape */ ],
    "genre":   [ /* same shape */ ],
    "frame_predictions": {
      "frame_hop_seconds": 1.0,
      "segment_probs": [[0.18, 0.13, 0.05, 0.11, 0.11, 0.41], ...],
      "arousal_probs": [[0.06, 0.27, 0.67], ...],
      "valence_probs": [[0.06, 0.40, 0.53], ...],
      "genre_probs":   [[0.05, 0.10, ...], ...],
      "arousal_reg":   [0.36, 0.38, 0.41, ...],
      "valence_reg":   [0.42, 0.45, 0.50, ...]
    },
    "audio_features": {
      "tempo_bpm": 123.0,
      "key": { "key": "F", "mode": "Minor", "confidence": 0.78 }
    }
  }
}
\end{lstlisting}

\section{CLI-команды для воспроизведения}\label{app:F}

\subsection{Подготовка данных}

\begin{lstlisting}[language=bash]
# Activate environment
source .venv/bin/activate          # Linux/Mac
.venv/Scripts/activate             # Windows

# DEAM
python scripts/prepare_deam.py --data-dir data/deam

# Harmonix (requires yt-dlp + ffmpeg)
python scripts/download_harmonix_audio.py --data-dir data/structure
python scripts/prepare_structure.py --data-dir data/structure

# GTZAN
python scripts/prepare_gtzan.py --data-dir data/gtzan
\end{lstlisting}

\subsection{Обучение моделей}

\begin{lstlisting}[language=bash]
# CNN v2 (default)
python scripts/train.py \
    --config configs/default.yaml \
    --deam-dir data/deam --structure-dir data/structure --gtzan-dir data/gtzan \
    --output-dir checkpoints/cnn_v2

# AST v2 (final)
python scripts/train.py \
    --config configs/ast.yaml \
    --deam-dir data/deam --structure-dir data/structure --gtzan-dir data/gtzan \
    --output-dir checkpoints/ast_v2

# CNN v3 (failed)
python scripts/train.py \
    --config configs/cnn_v3.yaml \
    --deam-dir data/deam --structure-dir data/structure --gtzan-dir data/gtzan \
    --output-dir checkpoints/cnn_v3
\end{lstlisting}

\subsection{Оценка моделей}

\begin{lstlisting}[language=bash]
python scripts/eval.py \
    --ckpt checkpoints/ast_v2/best.pt \
    --config configs/ast.yaml \
    --output results/eval_ast_v2.csv \
    --deam-dir data/deam --structure-dir data/structure --gtzan-dir data/gtzan
\end{lstlisting}

\subsection{Инференс на одном файле}

\begin{lstlisting}[language=bash]
python scripts/infer.py \
    --audio path/to/song.mp3 \
    --ckpt checkpoints/ast_v2/best.pt \
    --config configs/ast.yaml \
    --output results/timeline.json \
    --plot results/timeline.png
\end{lstlisting}

\subsection{Запуск веб-прототипа}

\begin{lstlisting}[language=bash]
# All three services with one command:
bun run scripts/dev-all.ts

# Or separately:
python scripts/serve_ml.py --ckpt checkpoints/ast_v2/best.pt --config configs/ast.yaml --port 8000
cd web/backend && bun run dev    # port 3000
cd web/frontend && bun run dev   # port 5173
\end{lstlisting}

Открыть в браузере: \url{http://localhost:5173}.

\subsection{Тесты и качество кода}

\begin{lstlisting}[language=bash]
pytest tests/ -v                   # 128 tests
ruff check musicml/ scripts/ tests/
\end{lstlisting}
